(a) A $\mathbb P^n$-bundle is a morphism $\pi:P\to X$ together with local $X$-isomorphisms $\pi^{-1}(U_i)\cong\mathbb P^n_{U_i}$ whose transitions, locally on each overlap, are induced by invertible linear substitutions in the homogeneous coordinates. Two such atlases are equivalent when their union has this property. The transitions satisfy the cocycle condition as projective automorphisms; matrices inducing them are determined only up to multiplication by a unit.

(b) A basis of $\mathcal E|_{U_i}$ identifies its symmetric algebra with $\mathcal O_{U_i}[x_0,\ldots,x_n]$. Changes of basis give invertible linear substitutions, so the corresponding Proj charts make $\mathbb P(\mathcal E)$ a $\mathbb P^n$-bundle.

(c) The assertion is immediate for $n=0$. We may otherwise work on one component of $X$, so $X$ is integral. In a nonempty trivializing open $U$, choose a relative hyperplane in $\mathbb P^n_U$. Its closure in $P$ is a Weil divisor. Since $P$ is regular, this divisor is Cartier and defines an invertible sheaf $\mathcal L$ extending $\mathcal O_{\mathbb P^n_U}(1)$.

On any connected trivializing open $V$, \exref{II.7.9}(a) gives $\mathcal L|_{\mathbb P^n_V}\cong\mathcal O_{\mathbb P^n_V}(d)\otimes\pi^*\mathcal M$. Restriction to the generic fibre gives $d=1$. After shrinking $V$ to trivialize $\mathcal M$, $\pi_*\mathcal L$ is free of rank $n+1$, and the evaluation map $\pi^*\pi_*\mathcal L\to\mathcal L$ is surjective. Thus $\mathcal E=\pi_*\mathcal L$ is locally free, and (7.12) gives $P\to\mathbb P(\mathcal E)$. On these opens it is the usual isomorphism $\mathbb P^n_V\to\mathbb P(\mathcal O_V^{n+1})$, so it is an isomorphism.

Regularity can be weakened to local factoriality. Polynomial rings over a factorial domain, and their localizations, are factorial, so $P$ is then locally factorial; the divisor extension and \exref{II.7.9}(a) apply unchanged.

(d) Parts (b) and (c) give surjectivity of $\mathcal E\mapsto\mathbb P(\mathcal E)$ onto the isomorphism classes of $\mathbb P^n$-bundles. By \exref{II.7.9}(b), two sheaves give the same bundle exactly when they differ by tensoring with an invertible sheaf. For $n=0$, all invertible sheaves are equivalent in this way and the only bundle is $X$.
